\newpage
\section{正定质量粒子因果场\(\dagger\)}
一般形式场写为：\begin{equation*}
    \Psi_\ell(x) =(\frac{1}{2\pi})^\frac{3}{2}\sum_{\sigma}\int d^3p \, \left\{\xi_\ell u_\ell(\boldsymbol{p},\sigma)a^{c\dagger}(\boldsymbol{k},\sigma) e^{-ipx} + \lambda_\ell v_\ell(\boldsymbol{p},\sigma) a(\boldsymbol{k},\sigma)e^{ipx} \right\}
\end{equation*}
若理论本身无荷，就有$a^{\dagger}(\boldsymbol{k},\sigma)=a^{c\dagger}(\boldsymbol{k},\sigma)=a^{\dagger}(\boldsymbol{k},\sigma)$；系数$u_\ell(\boldsymbol{p},\sigma)$和$v_\ell(\boldsymbol{p},\sigma)$的变换是
\begin{equation*}
    \begin{aligned}
            &\sum_{\bar{\sigma}}
            u_{\ell}(\boldsymbol{p},\bar{\sigma})
            \sqrt{\frac{p^0}{(\Lambda p)^0}} 
            D_{\sigma \bar{\sigma}}^{j_n}
            \!\left\{W^{-1}(\Lambda,p)\right\} =
            \sum_{\bar{\ell}}
            \mathscr{D}_{\ell\bar{\ell}}(\Lambda^{-1})
            u_{\bar{\ell}}(\Lambda \boldsymbol{p},\sigma)\\
            &\sum_{\bar{\sigma}}
            v_{\ell}(\boldsymbol{p},\bar{\sigma})
            \sqrt{\frac{p^0}{(\Lambda p)^0}} 
            D_{\sigma \bar{\sigma}}^{j_n*}
            \!\left\{W^{-1}(\Lambda,p)\right\} =
            \sum_{\bar{\ell}}
            \mathscr{D}_{\ell\bar{\ell}}(\Lambda^{-1})
            v_{\bar{\ell}}(\Lambda \boldsymbol{p},\sigma)
            \end{aligned}
\end{equation*}
\begin{center}
    \vspace{2ex}
    \large\textbf{***}
    \vspace{2ex}
\end{center}
通过 $SU(2)_L \times SU(2)_R$ 代数构建Lorentz群的有限维表示后。指标$\ell$在这里被替换为一对指标$a$、$b$：所以一般形式场可以写为：\begin{equation}
    \Psi_{ab}(x) =(\frac{1}{2\pi})^\frac{3}{2}\sum_{\sigma}\int d^3p \, \left\{\xi_{ab} u_{ab}(\boldsymbol{p},\sigma)a^{c\dagger}(\boldsymbol{k},\sigma) e^{-ipx} + \lambda_{ab} v_{ab}(\boldsymbol{p},\sigma) a(\boldsymbol{k},\sigma)e^{ipx} \right\}
\end{equation}
以及系数函数的选择变为：
\begin{subequations}
    \begin{align}
            &\sum_{\bar{\sigma}}
            u_{ab}(\boldsymbol{p},\bar{\sigma})
            \sqrt{\frac{p^0}{(\Lambda p)^0}} 
            D_{\sigma \bar{\sigma}}^{j_n}
            \!\left\{W^{-1}(\Lambda,p)\right\} =
            \sum_{\bar{a}\bar{b}}
            \mathscr{D}_{{ab};\bar{a}\bar{{b}}}(\Lambda^{-1})
            u_{\bar{a}\bar{b}}(\Lambda \boldsymbol{p},\sigma)
            \\
             &\sum_{\bar{\sigma}}
            v_{ab}(\boldsymbol{p},\bar{\sigma})
            \sqrt{\frac{p^0}{(\Lambda p)^0}} 
            D_{\sigma \bar{\sigma}}^{j_n*}
            \!\left\{W^{-1}(\Lambda,p)\right\} =
            \sum_{\bar{a}\bar{b}}
            \mathscr{D}_{{ab};\bar{a}\bar{{b}}}(\Lambda^{-1})
            v_{\bar{a}\bar{b}}(\Lambda \boldsymbol{p},\sigma)
            \end{align}
\end{subequations}
首先确定系数的选择；考虑剩下的两组Poincar\'e变换，依次考察$\operatorname{Boost}$和$\operatorname{Rotation}$；
\vspace{2ex}
\begin{center}
    \large\textbf{boost}  
\end{center}
\vspace{2ex}
令$p^\mu=(p^0,\boldsymbol{p})$，也就是任意动量系；选择$\Lambda=L^{-1}(p)$是任意系到静止系标准Lorentz变换，$\Lambda p=k=(m,0,0,0)$，所以：
\[\begin{aligned}
    &W(\Lambda,p)=1\to\sum_{\bar{\sigma}}D_{\sigma \bar{\sigma}}^{j_n}\!\left\{W^{-1}(\Lambda,p)\right\} =1\\
    &\qquad(\Lambda p)^0\to m^0
\end{aligned}\]
这种情况下，系数函数表述为：
\begin{subequations}
    \begin{align}\label{eq:8.38a}
       &\sum_{\bar{\sigma}}u_{ab}(\boldsymbol{p},\bar{\sigma})\delta_{\sigma\bar{\sigma}}\sqrt{\frac{p^0}{m}}
        =\sum_{\bar{a}\bar{b}}\mathscr{D}_{ab;\bar{a}\bar{b}}(L(p))u_{\bar{a}\bar{b}}(0,\sigma)\\
       &\sum_{\bar{\sigma}}v_{ab}(\boldsymbol{p},\bar{\sigma})\delta_{\sigma\bar{\sigma}}\sqrt{\frac{p^0}{m}}
        =\sum_{\bar{a}\bar{b}}\mathscr{D}_{ab;\bar{a}\bar{b}}(L())v_{\bar{a}\bar{b}}(0,\sigma)
\end{align}
\end{subequations}
\vspace{2ex}
\begin{center}
   \large \textbf{rotation}  
\end{center}
\vspace{2ex}
令$p^\mu=(p^0,0,0,0)$，也就是静止系；选择$\Lambda=R(\theta)$也就是在静止系下保持$p^\mu=(p^0,0,0,0)$不变的旋转，那么有$R(p)=p$以及$\Lambda p=p$，所以：
\[W(\Lambda,p)=R\]
\color{blue}由于$\boldsymbol{J}=\boldsymbol{A}+\boldsymbol{B}$，$Winger$的转动表征为总的自旋$\boldsymbol{J}$,
$$ D_{\sigma\bar{\sigma}}^{(j)}(R^{-1}) = \left[ \exp\left( +i \boldsymbol{\theta}\cdot\boldsymbol{J}^{(j)} \right) \right]_{\sigma\bar{\sigma}} $$
场的自旋拆为$\boldsymbol{A}+\boldsymbol{B}$
\[M_{{ab};\bar{a}\bar{{b}}}(R^{-1})=\left[\exp\left( +i \boldsymbol{\theta}\cdot\boldsymbol{A}+i \boldsymbol{\theta}\cdot\boldsymbol{B}\right) \right]_{{ab};\bar{a}\bar{{b}}}\]\normalcolor
在这种分解下系数函数的选择\footnote{这个方程也叫做交织方程}变为：
\begin{subequations}\label{u_0}
    \begin{align}
            &\sum_{\bar{\sigma}} 
            \left[ \exp\left( -i \boldsymbol{\theta}\cdot\boldsymbol{J}^{(j)} \right) \right]_{\sigma\bar{\sigma}}u_{ab}(0,\bar{\sigma}) =
            \sum_{\bar{a}\bar{b}}
            \left[\exp\left(  +i \boldsymbol{\theta}\cdot\boldsymbol{A}+i \boldsymbol{\theta}\cdot\boldsymbol{B} \right) \right]_{{ab};\bar{a}\bar{{b}}}
            u_{\bar{a}\bar{b}}(0,\sigma)
            \\
            &\sum_{\bar{\sigma}} 
            \left[ \exp\left( +i \boldsymbol{\theta}\cdot\boldsymbol{J}^{(j)} \right) \right]_{\sigma\bar{\sigma}}v_{ab}(0,\bar{\sigma}) =
            \sum_{\bar{a}\bar{b}}
            \left[\exp\left(  +i \boldsymbol{\theta}\cdot\boldsymbol{A}+i \boldsymbol{\theta}\cdot\boldsymbol{B} \right) \right]_{{ab};\bar{a}\bar{{b}}}
            v_{\bar{a}\bar{b}}(0,\sigma)
            \end{align}
\end{subequations}
\color{blue}下面只需要在静止系中确定系数的形式；之后通过选择一个标准的$boost$来获得动量。\normalcolor
\begin{center}
    \vspace{2ex}
    \large\text{系数$u(0,\sigma)$和$v(0,\sigma)$和角动量加法}
\end{center}
\vspace{2ex}
交织方程(\refeq{u_1})在$\boldsymbol{\theta}$ 在 $\boldsymbol{\theta}=0$ 处展开至一阶对比两边 $\boldsymbol{\theta}$ 的系数，，到静止系下系数函数的代数关系变为
\begin{subequations}
    \begin{align}
       -\sum_{\bar{\sigma}} u_{ab}(0,\bar{\sigma}) \left[ \boldsymbol{J}^{(j)} \right]_{\bar{\sigma}\sigma} &= 
       \sum_{\bar{a}\bar{b}} \left[ \mathbf{A}_{a\bar{a}} \delta_{b\bar{b}} + \delta_{a\bar{a}} \mathbf{B}_{b\bar{b}} \right] u_{\bar{a}\bar{b}}(0,\sigma)
        \\
        +\sum_{\bar{\sigma}} v_{ab}(0,\bar{\sigma}) \left[ \boldsymbol{J}^{(j)} \right]_{\bar{\sigma}\sigma} &= 
        \sum_{\bar{a}\bar{b}} \left[ \mathbf{A}_{a\bar{a}} \delta_{b\bar{b}} + \delta_{a\bar{a}} \mathbf{B}_{b\bar{b}} \right] v_{\bar{a}\bar{b}}(0,\sigma)
    \end{align}
\end{subequations}
\color{blue}这是一个极其深刻的物理方程。这是两个角动量算符的直积之和：$\boldsymbol{J}^{(A)} \otimes \mathbb{I} + \mathbb{I} \otimes \boldsymbol{J}^{(B)}$；而等式左边则是单一的总角动量算符 $\boldsymbol{J}^{(j)}$。实际上,系数 $u_{ab}(0, \sigma)$ 的等式，本质上就是角动量相加方法\normalcolor。为了清晰地展现这一点，首先明确系统的两种物理表象：
\begin{enumerate}
    \item 非耦合表象：
    场的洛伦兹张量指标 $(a,b)$ 属于表示空间 $V_A \otimes V_B$。如果我们将 $A$ 和 $B$ 视为两个独立的子系统，，对应的正交归一基底为：
    \begin{equation}
        |A,a; B,b\rangle \equiv |A,a\rangle \otimes |B,b\rangle, \quad a = -A, \dots, A, \quad b = -B, \dots, B.
    \end{equation}
    这组基底构成了系统的完备集，满足 $\sum_{a,b} |A,a; B,b\rangle \langle A,a; B,b| = \mathbb{I}$。

    \item \textbf{耦合表象}：
    对于静止参考系下的物理单粒子态，必须在旋转下按确定自旋 $j$ 的不可约表示变换。我们用总角动量基底来标记物理态：
    \begin{equation}
        |j, \bar{\sigma}\rangle, \quad \bar{\sigma} = -j, \dots, j.
    \end{equation}
\end{enumerate}
在耦合表象中插入非耦合表象的完备基底，就可以表达出$|j, \bar{\sigma}\rangle$；
\begin{equation}
    |j, \bar{\sigma}\rangle = \mathbb{I} \, |j, \bar{\sigma}\rangle = \sum_{a,b} |A,a; B,b\rangle \langle A,a; B,b | j, \bar{\sigma}\rangle.
    \label{eq:QM_expansion}
\end{equation}
这里的展开系数 $\langle A,a; B,b | j, \bar{\sigma}\rangle$，，定义就是连接这两个空间的 \textbf{Clebsch-Gordan 系数}
\vspace{2ex}\\
系数 $u_{ab}(0, \bar{\sigma})$ 的物理意义;是将物理自旋态“投影”到具体的场张量分量上。不严谨但形象的说，$u_{ab}(0, \bar{\sigma})$ 充当了态矢量的展开系数：
\begin{equation}
    |j, \bar{\sigma}\rangle \sim \sum_{a,b} |A,a; B,b\rangle \, u_{ab}(0, \bar{\sigma}).
    \label{eq:Field_expansion}
\end{equation}
比较式 \eqref{eq:QM_expansion} 与式 \eqref{eq:Field_expansion}，结合我们刚刚推导出的交织方程：方程要求 $u_{ab}(0, \bar{\sigma})$ 必须满足总角动量守恒的变换规则，与 C-G 系数所满足的代数方程是完全一致的。
\vspace{2ex}\\
从群表示论的严格角度来看， $SU(2)$ 的角动量加法 
\begin{equation}
V^{(A)} \otimes V^{(B)}
=
\bigoplus_{j=|A-B|}^{A+B}
V^{(j)} .
\label{eq:CG_decomposition}
\end{equation} 中，一个自旋 $j$ 的不可约表示在直和分解中只出现一次。这意味着：够满足这种旋转协变性的基底变换矩阵是\textbf{唯一确定的}（Schur 引理的推论）。
这种唯一性说明了 $u_{ab}(0, \bar{\sigma})$是正比于C-G 系数：
\begin{equation}
    u_{ab}(0, \bar{\sigma}) = N \langle A,a; B,b \,|\, j, \bar{\sigma}\rangle.
    \label{eq:uCG_proportional}
\end{equation}
这里 $N$ 是一个与群论旋转无关的整体常数。约定取 $N = \frac{1}{\sqrt{2m}}$，，终得到
\begin{equation}
    u_{ab}(0, \bar{\sigma}) = \frac{1}{\sqrt{2m}} \langle A,a; B,b \,|\, j, \bar{\sigma}\rangle.
\end{equation}
利用角动量态在反演下的相因子变换性质，及 C-G 系数的内在对称性，可以推导出 $v_{ab}$ 对应于物理自旋方向相反的态，带有一个标准的 Condon-Shortley 相位：
\begin{equation}
    v_{ab}(0,\bar{\sigma}) = \xi^{AB} \, (-1)^{j-\bar{\sigma}} \langle A, a ; B, b \,|\, j, -\bar{\sigma} \rangle
    \label{eq:v_CG}
\end{equation}
其中 $\xi^{AB}$ 是与粒子的内禀宇称和场定义相关的纯相因子（通常取为 $\pm 1$ 或 $\pm i$）

\vspace{2ex}
\begin{center}
    \large\text{标准$boost$获得动量}
\end{center}
\vspace{2ex}
已经确定了静止参考系下的系数 $u_{ab}(0, \sigma)$ 正比于 C-G 系数。那么如何获得具有任意有限动量 $\boldsymbol{p}$ 的物理态系数 $u_{ab}(\boldsymbol{p}, \sigma)$？ 

采取的办法是做一个$boost$获得有限动，由于 $\mathscr{D}(\Lambda)$ 是洛伦兹群的表示矩阵，满足 $\mathscr{D}(L(p)) \mathscr{D}(L^{-1}(p)) = \mathbb{I}$。我们在方程\eqref{eq:8.38a} 两边同时左乘 $\mathscr{D}(L(p))$，以解出动量为 $\boldsymbol{p}$ 时的场系数：
\begin{equation}
    \sum_{\bar{\sigma}}u_{ab}(\boldsymbol{p},\bar{\sigma})\delta_{\sigma\bar{\sigma}}
        =\sqrt{\frac{m}{p^0}}\sum_{\bar{a}\bar{b}}\mathscr{D}_{ab;\bar{a}\bar{b}}(L(p))u_{\bar{a}\bar{b}}(0,\sigma)
    \label{eq:u_boosted}
\end{equation}
物理图像到这里变得无比清晰：\color{blue}要构造一个动量为 $\boldsymbol{p}$、自旋投影为 $\sigma$ 的系数函数，我们只需要先在静止系用 C-G 系数把它“拼装”出来，然后再用一个标准的洛伦兹 Boost 矩阵 $\mathscr{D}(L(p))$ 把整个张量“加速”到动量 $p$ 即可。\normalcolor

为了让这个公式具备实际的可计算性，需要写出 $\mathscr{D}(L(p))$ 在 $(A,B)$ 表示下的具体指数形式。设 Boost 的快度矢量为 $\boldsymbol{\phi} = \hat{\boldsymbol{p}} \text{arctanh}(|\boldsymbol{p}|/p^0)$\footnote{
为了计算有限动量态的旋量系数，引入了“Rapidity” $\phi$ 作为核心参数，其定义为：
\begin{align*}
    \cosh\phi &= \frac{\sqrt{\boldsymbol{k}^2+m^2}}{m} \\
    \sinh\phi&= \frac{|\boldsymbol{k}|}{m}
\end{align*}
这种参量化的最大物理优势在于，同向的 Boost 变换具有简单的加法性质：$L(\bar{\phi})L(\phi) = L(\bar{\phi}+\phi)$。
}对应的生成元为 $\boldsymbol{K}$。在 $(A,B)$ 表示中，生成元满足：
\begin{equation}
    \boldsymbol{J} = \mathbf{A} + \mathbf{B}, \quad \boldsymbol{k} = -i\mathbf{A} + i\mathbf{B}
\end{equation}
所以，有限 Boost 操作的表示矩阵可以彻底分解为左右两部分的张量积：
\begin{align}
    \mathscr{D}\bigl(L(q)\bigr) &= \exp\left( -i \boldsymbol{\phi} \cdot \boldsymbol{k} \right) \nonumber \\
    &= \exp\left( -i \boldsymbol{\phi} \cdot (-i\mathbf{A} + i\mathbf{B}) \right) \nonumber \\
    &= \exp\left( -\boldsymbol{\phi} \cdot \mathbf{A} \right) \otimes \exp\left( +\boldsymbol{\phi} \cdot \mathbf{B} \right)
\end{align}
将此矩阵代入 \eqref{eq:u_boosted}，，结合静止系下的 C-G 系数结果，最终得到了场论中\textbf{最一般的动力学系数展开式}：
\begin{equation}
   \sum_{\bar{\sigma}} u_{ab}(\boldsymbol{p}, \bar{\sigma}) = \frac{1}{\sqrt{2p^0}} \sum_{\bar{a}\bar{b}} \left[ e^{-\boldsymbol{\phi} \cdot \mathbf{A}} \right]_{a\bar{a}} \left[ e^{+\boldsymbol{\phi} \cdot \mathbf{B}} \right]_{b\bar{b}} \langle A,a; B,b \,|\, j,\bar{\sigma} \rangle
\end{equation}
而反粒子对应的系数的是：
\begin{equation}
    v_{ab}(\boldsymbol{p}, \sigma)=(-)^{j+\sigma}u_{ab}(\boldsymbol{p}, -\sigma)
\end{equation}

\vspace{2ex}
\begin{center}
    \large\text{微观因果性和统计性质}
\end{center}
\vspace{2ex}



\subsection{因果标量场 }
一般形式的场是：
\begin{equation}
    \Psi_{ab}(x) =(\frac{1}{2\pi})^\frac{3}{2}\sum_{\sigma}\int d^3p \, \left\{ u_{ab}(\boldsymbol{q},\sigma)b^\dagger(\boldsymbol{q},\sigma) e^{-ipx} + (-)^{2B}(-)^{j+\sigma}u_{ab}(\boldsymbol{q},\sigma) a(\boldsymbol{q},\sigma)e^{ipx} \right\}
\end{equation}
上述我们知道静止系系数 $u(0)$ 和 $v(0)$ 是正比于Clebsch-Gordan  系数 
\begin{equation}
    u_{ab}(0, \sigma) = \frac{1}{\sqrt{2m}} \langle A, a; B, b | j, \sigma \rangle
    \label{eq:u0_general}
\end{equation}
通过 Boost 获得有限动量系数 $u(\boldsymbol{q})$ 和 $v(\boldsymbol{q})$的一般公式是：
\begin{equation}
    u_{ab}(\boldsymbol{q}, \sigma) = \frac{1}{\sqrt{2q^0}} \sum_{\bar{a},\bar{b}} \left[e^{-\boldsymbol{\phi}\cdot\mathbf{J}^{(A)}}\right]_{a\bar{a}} \left[e^{+\boldsymbol{\phi}\cdot\mathbf{J}^{(B)}}\right]_{b\bar{b}} \langle A,\bar{a}; B,\bar{b} | j, \sigma \rangle
    \label{eq:u_finite_general}
\end{equation}
其中$\boldsymbol{\phi} = \hat{\boldsymbol{q}} \text{arctanh}(|\boldsymbol{q}|/q^0)$
\vspace{0.5cm}
\begin{center}
    \large\text{$(A, B) = (0, 0)$}
\end{center}
\vspace{0.5cm}
对于标量场，$(A, B) = (0, 0)$，粒子自旋 $j = 0$。此时：
\begin{itemize}
    \item 角动量生成元全为零：$\mathbf{J}^{(A)} = \mathbf{J}^{(B)} = 0$。
    \item 表示指标只有一个值：$a = 0$，$b = 0$。
    \item 自旋投影只有一个值：$\sigma = 0$。
    \item $j=0$根据自旋统计表明是波色子。
\end{itemize}
代入式 \eqref{eq:u0_general}，唯一的 C-G 系数为 $\langle 0, 0; 0, 0 | 0, 0 \rangle = 1$。因此静止系下的系数为：
\begin{equation}
    u(0, 0) = \frac{1}{\sqrt{2m}}
\end{equation}
\vspace{2ex}
对于标量场，$\mathbf{J}^{(A)} = 0$ 且 $\mathbf{J}^{(B)} = 0$。因此，指数演化矩阵完全退化为恒等映射：
\begin{equation}
    \exp(-\boldsymbol{\phi}\cdot\mathbf{0}) = 1, \quad \exp(+\boldsymbol{\phi}\cdot\mathbf{0}) = 1
\end{equation}
代入式 \eqref{eq:u_finite_general}，并利用 $\langle 0,0;0,0|0,0\rangle = 1$，有限动量系数极为简单：
\begin{equation}
    u(\boldsymbol{q}, 0) = \frac{1}{\sqrt{2q^0}}
\end{equation}
所以最终因果标量场是：
\begin{equation}
    \phi(x) =(\frac{1}{2\pi})^\frac{3}{2}\int \frac{d^3p }{\sqrt{2q^0}}\, \left\{ b^\dagger(\boldsymbol{q},0) e^{-ipx} + a(\boldsymbol{q},0)e^{ipx} \right\}
\end{equation}若粒子与反粒子不同（带电标量场，$a \neq b$），系统具有整体 $U(1)$ 规范对称性，对应守恒荷算符：
\begin{equation}
    Q = \int d^3p\left(a^\dagger(\boldsymbol{q})a(\boldsymbol{q}) - b^\dagger(\boldsymbol{q})b(\boldsymbol{q})\right)
\end{equation}
利用玻色对易关系验证：
\begin{equation}
    [Q, \phi(x)] = -\phi(x), \quad [Q, \phi^\dagger(x)] = +\phi^\dagger(x)
\end{equation}
$\phi$ 每作用一次使系统总电荷减少 $1$（湮没正粒子或产生反粒子），$\phi^\dagger$ 使总电荷增加 $1$。
为了后续计算传播子，需要考察系数函数的完备性。对于标量场，自旋 $j = 0$，只有一个极化态 $\sigma = 0$：
\begin{equation}
    \sum_{\sigma=0} u(\boldsymbol{q},\sigma)\, u^*(\boldsymbol{q},\sigma) = \frac{1}{2q^0}, \quad \sum_{\sigma=0} v(\boldsymbol{q},\sigma)\, v^*(\boldsymbol{q},\sigma) = \frac{1}{2q^0}
\end{equation}


\subsection{Weyl 场}
一般形式的场是：
\begin{equation}
    \Psi_{ab}(x) =(\frac{1}{2\pi})^\frac{3}{2}\sum_{\sigma}\int d^3p \, \left\{ u_{ab}(\boldsymbol{q},\sigma)b^\dagger(\boldsymbol{q},\sigma) e^{-ipx} + (-)^{2B}(-)^{j+\sigma}u_{ab}(\boldsymbol{q},\sigma) a(\boldsymbol{q},\sigma)e^{ipx} \right\}
\end{equation}
上述我们知道静止系系数 $u(0)$ 和 $v(0)$ 是正比于Clebsch-Gordan  系数 
\begin{equation}
    u_{ab}(0, \sigma) = \frac{1}{\sqrt{2m}} \langle A, a; B, b | j, \sigma \rangle
    \label{eq:u0_general_1/2}
\end{equation}
通过 Boost 获得有限动量系数 $u(\boldsymbol{q})$ 和 $v(\boldsymbol{q})$的一般公式是：
\begin{equation}
    u_{ab}(\boldsymbol{q}, \sigma) = \frac{1}{\sqrt{2q^0}} \sum_{\bar{a},\bar{b}} \left[e^{-\boldsymbol{\phi}\cdot\mathbf{J}^{(A)}}\right]_{a\bar{a}} \left[e^{+\boldsymbol{\phi}\cdot\mathbf{J}^{(B)}}\right]_{b\bar{b}} \langle A,\bar{a}; B,\bar{b} | j, \sigma \rangle
    \label{eq:u_finite_general_1/2}
\end{equation}
其中$\boldsymbol{\phi} = \hat{\boldsymbol{q}} \text{arctanh}(|\boldsymbol{q}|/q^0)$
\vspace{0.5cm}
\begin{center}
    \large\text{$(A, B) = (1/2, 0)$或$(1/2, 0)$}
\end{center}
\vspace{0.5cm}
对于标量场，$(A, B) = (0, 0)$或$(1/2, 0)$，粒子自旋 $j = 1/2$。此时：
\begin{itemize}
    \item 角动量生成元全为零：$\mathbf{J}^{(A)} = 1/2$或${J}^{(B)} = 1/2$
    \item 表示指标只有一个值：$a = 1/2$或$b = 1/2$。
    \item 自旋投影只有一个值：$\sigma = 1/2$。
    \item $j=1/2$根据自旋统计表明是费米子。
\end{itemize}
代入式 \eqref{eq:u0_general_1/2}，唯一的 C-G 系数为 $\langle 1/2, a; 0, 0 | 1/2, \sigma \rangle = \delta_{a\sigma}$。因此静止系下的系数为：
\begin{equation}
    u_a(0, \sigma) = \frac{1}{\sqrt{2m}} \langle \tfrac{1}{2}, a;\, 0, 0\, |\, \tfrac{1}{2}, \sigma \rangle = \frac{1}{\sqrt{2m}}\, \delta_{a\sigma}
\end{equation}
\textbf{静止系系数的具体列向量形式}将 $a = +1/2, -1/2$ 按顺序排列为列向量的两个分量，我们显式写出两个自旋投影对应的系数：

\textbf{对于 $\sigma = +1/2$：} $u_a(0, +1/2) = \frac{1}{\sqrt{2m}}\delta_{a, +1/2}$，即
\begin{equation}
    u(0, +\tfrac{1}{2}) = \frac{1}{\sqrt{2m}} \begin{pmatrix} 1 \\ 0 \end{pmatrix}=\frac{1}{\sqrt{2m}}\, \xi_{\frac{1}{2}}
\end{equation}

\textbf{对于 $\sigma = -1/2$：} $u_a(0, -1/2) = \frac{1}{\sqrt{2m}}\delta_{a, -1/2}$，即
\begin{equation}
    u(0, -\tfrac{1}{2}) = \frac{1}{\sqrt{2m}} \begin{pmatrix} 0 \\ 1 \end{pmatrix}=\frac{1}{\sqrt{2m}}\, \xi_{-\frac{1}{2}}
\end{equation}

其中 $\xi_{+1/2} = \binom{1}{0}$ 和 $\xi_{-1/2} = \binom{0}{1}$ 是标准的二分量自旋基矢。
\vspace{2ex}
对于 $(A, B) = (1/2, 0)$ ，Boost 生成元为 $\mathbf{K} = -i\mathbf{J}^{(A)} + i\mathbf{J}^{(B)} = -i\frac{\boldsymbol{\sigma}}{2}$。因此，一般 Boost 分解公式简化为：
\begin{equation}
    \mathscr{D}_{1/2,0}(L(p)) = \exp\left(-\boldsymbol{\phi}\cdot\mathbf{J}^{(A)}\right) \otimes 1 = \exp\left(-\frac{\boldsymbol{\phi}\cdot\boldsymbol{\sigma}}{2}\right)
\end{equation}
利用泡利矩阵的性质 $(\hat{\boldsymbol{q}}\cdot\boldsymbol{\sigma})^2 = I$，指数矩阵可以精确展开为：
\begin{equation}
    \exp\left(-\frac{\phi}{2}\, \hat{\boldsymbol{q}}\cdot\boldsymbol{\sigma}\right) = \cosh\frac{\phi}{2}\cdot I - \sinh\frac{\phi}{2}\cdot \hat{\boldsymbol{q}}\cdot\boldsymbol{\sigma}
\end{equation}
利用双曲半角公式（其中 $\cosh\phi = q^0/m$，$\sinh\phi = |\boldsymbol{q}|/m$）：
\begin{align}
    \cosh\frac{\phi}{2} &= \sqrt{\frac{\cosh\phi + 1}{2}} = \sqrt{\frac{q^0 + m}{2m}} \\
    \sinh\frac{\phi}{2} &= \sqrt{\frac{\cosh\phi - 1}{2}} = \sqrt{\frac{q^0 - m}{2m}}
\end{align}
代入得：
\begin{equation}
    \exp\left(-\frac{\boldsymbol{\phi}\cdot\boldsymbol{\sigma}}{2}\right) = \frac{1}{\sqrt{2m}}\left[\sqrt{q^0 + m}\cdot I - \sqrt{q^0 - m}\cdot \hat{\boldsymbol{q}}\cdot\boldsymbol{\sigma}\right]
    \label{eq:weyl_boost_expanded}
\end{equation}
带入\eqref{eq:u_finite_general_1/2}
\begin{equation}
   \begin{aligned}
     u(\boldsymbol{q}, \sigma) &= \sqrt{\frac{1}{2q^0}} \cdot \frac{1}{\sqrt{2m}} \cdot \frac{\sqrt{q^0+m}\cdot I - \sqrt{q^0-m}\cdot\hat{\boldsymbol{q}}\cdot\boldsymbol{\sigma}}{\sqrt{2m}} \xi_\sigma\\
    & = \frac{1}{2\sqrt{mq^0}} \frac{\sqrt{q^0+m}\cdot I - \sqrt{q^0-m}\cdot\hat{\boldsymbol{q}}\cdot\boldsymbol{\sigma}}{\sqrt{2m}}\xi_\sigma
    \label{eq:weyl_u_p}
   \end{aligned}
\end{equation}
将式 \eqref{eq:weyl_u_p}  中的 Boost 矩阵的平方写成四维协变形式。注意到：
\begin{equation}
    \exp\left(-\boldsymbol{\phi}\cdot\boldsymbol{\sigma}\right) = \left[\exp\left(-\frac{\boldsymbol{\phi}\cdot\boldsymbol{\sigma}}{2}\right)\right]^2 = \cosh\phi\cdot I - \sinh\phi\cdot\hat{\boldsymbol{q}}\cdot\boldsymbol{\sigma} = \frac{q^0 I - \boldsymbol{q}\cdot\boldsymbol{\sigma}}{m}
\end{equation}
定义 $\sigma^\mu = (I, \boldsymbol{\sigma})$，则上式可紧凑地写为：
\begin{equation}
    \exp(-\boldsymbol{\phi}\cdot\boldsymbol{\sigma}) = \frac{p_\mu\sigma^\mu}{m}
    \label{eq:boost_squared}
\end{equation}
这一恒等式是联系 Boost 矩阵与四维动量的桥梁，将在极化求和中发挥关键作用。计算旋量系数的完备性关系。利用式 \eqref{eq:weyl_u_p}，对两个自旋投影 $\sigma = \pm 1/2$ 求和：
\begin{equation}
    \sum_\sigma u(\boldsymbol{q},\sigma)\, u^\dagger(\boldsymbol{q},\sigma) = \frac{m}{q^0}\, \exp\left(-\frac{\boldsymbol{\phi}\cdot\boldsymbol{\sigma}}{2}\right) \underbrace{\left(\frac{1}{2m}\sum_\sigma \xi_\sigma\xi_\sigma^\dagger\right)}_{= I/(2m)} \exp\left(-\frac{\boldsymbol{\phi}\cdot\boldsymbol{\sigma}}{2}\right)^\dagger
\end{equation}
由于 $\sum_\sigma \xi_\sigma\xi_\sigma^\dagger = I$（二分量基矢的完备性），以及 $\exp(-\frac{\boldsymbol{\phi}\cdot\boldsymbol{\sigma}}{2})$ 是厄米矩阵（泡利矩阵厄米），上式化为：
\begin{equation}
    \sum_\sigma u(\boldsymbol{q},\sigma)\, u^\dagger(\boldsymbol{q},\sigma) = \frac{1}{2q^0}\, \exp(-\boldsymbol{\phi}\cdot\boldsymbol{\sigma})
\end{equation}
代入四维协变形式 \eqref{eq:boost_squared}：
\begin{equation}
    \boxed{\sum_\sigma u(\boldsymbol{q},\sigma)\, u^\dagger(\boldsymbol{q},\sigma) = \frac{p_\mu\sigma^\mu}{2mq^0}}
    \label{eq:weyl_pol_u}
\end{equation}

对于 $v$ 系数，由于 $\sum_\sigma \eta_\sigma\eta_\sigma^\dagger = (i\sigma_2)I(i\sigma_2)^\dagger = I$，极化求和与 $u$ 完全相同：
\begin{equation}
    \boxed{\sum_\sigma v(\boldsymbol{q},\sigma)\, v^\dagger(\boldsymbol{q},\sigma) = \frac{p_\mu\sigma^\mu}{2mq^0}}
    \label{eq:weyl_pol_v}
\end{equation}
与标量场的标量极化求和 $1/(2q^0)$ 对比，Weyl 场的极化求和是一个 $2 \times 2$ 矩阵 $p_\mu\sigma^\mu/(2mq^0)$，其中包含了非平凡的自旋信息。正是这个矩阵中的动量 $p_\mu$，在下一步的因果性验证中将被替换为微商，从而带来决定性的符号翻转。



\paragraph{分立对称性：宇称破缺}

空间反演使 $\boldsymbol{q} \to -\boldsymbol{q}$，即快度 $\boldsymbol{\phi} \to -\boldsymbol{\phi}$。在 $(1/2, 0)$ 表示中，Boost 矩阵在宇称下的行为为：
\begin{equation}
    \mathcal{P}: \quad \underbrace{\exp\left(-\frac{\boldsymbol{\phi}\cdot\boldsymbol{\sigma}}{2}\right)}_{\text{左手 }(1/2,0) \text{ 的 Boost}} \longrightarrow \underbrace{\exp\left(+\frac{\boldsymbol{\phi}\cdot\boldsymbol{\sigma}}{2}\right)}_{\text{右手 }(0,1/2) \text{ 的 Boost}}
\end{equation}
这表明宇称操作将 $(1/2, 0)$ 的 Boost 矩阵变成了 $(0, 1/2)$ 的 Boost 矩阵，两者符号相反、不可等同。

因此，\textbf{单独的左手 Weyl 场在宇称变换下无法封闭回到自身}。如果物理理论要保持宇称对称性（如电磁力、强力），就必须同时引入左手 $(1/2, 0)$ 和右手 $(0, 1/2)$，构造它们的直和 $(1/2, 0) \oplus (0, 1/2)$。这就是下一节狄拉克形式体系的根本动机。

\subsection{狄拉克形式体系 }
两个 Weyl 场，它们各自只包含一种手征性：
\[
\Psi_L\in\Bigl(\tfrac12,0\Bigr),\qquad
\Psi_R\in\Bigl(0,\tfrac12\Bigr).
\]
在 Weyl 场中，$(1/2, 0)$ 表示的生成元为：
\begin{equation}
    \boldsymbol{A} = \frac{1}{2}(\boldsymbol{J} + i\boldsymbol{K}) = \frac{\boldsymbol{\sigma}}{2}, \quad \boldsymbol{B} = \frac{1}{2}(\boldsymbol{J} - i\boldsymbol{K}) = 0
\end{equation}
空间反演使 $\boldsymbol{K} \to -\boldsymbol{K}$，$\boldsymbol{J} \to \boldsymbol{J}$，因此：
\begin{equation}
    \mathcal{P}: \quad \boldsymbol{A} = \frac{1}{2}(\boldsymbol{J} + i\boldsymbol{K}) \longrightarrow \frac{1}{2}(\boldsymbol{J} - i\boldsymbol{K}) = \boldsymbol{B}
\end{equation}
这表明宇称操作将 $(1/2, 0)$ 表示映射为 $(0, 1/2)$ 表示。
为了在宇称变换下封闭，我们必须将两者组合为直和表示：
\begin{equation}
    \left(\frac{1}{2}, 0\right) \oplus \left(0, \frac{1}{2}\right)
\end{equation}
这便是四分量狄拉克旋量的来源。将四分量 Dirac 旋量 $\psi$ 构造为左右手 Weyl 旋量的列向量：
\begin{equation}
    \psi = \begin{pmatrix} \psi_R \\ \psi_L \end{pmatrix}, \quad \psi_R \in \left(0, \frac{1}{2}\right), \quad \psi_L \in \left(\frac{1}{2}, 0\right)
\end{equation}
这种将左右手分量明确分开的写法称为手征表示（Chiral/Weyl Representation）。

在洛伦兹变换下，上下两个二分量旋量按各自的表示独立变换。因此，直和表示的 $4 \times 4$ 变换矩阵是块对角化的：
\begin{equation}
    \mathscr{D}^{(1/2)}[\Lambda] = \begin{pmatrix} D^{(0,1/2)}[\Lambda] & 0 \\ 0 & D^{(1/2,0)}[\Lambda] \end{pmatrix}
\end{equation}
其中 $D^{(1/2,0)}[\Lambda] = e^{-\frac{1}{2}\boldsymbol{\phi}\cdot\boldsymbol{\sigma}}$ 和 $D^{(0,1/2)}[\Lambda] = e^{+\frac{1}{2}\boldsymbol{\phi}\cdot\boldsymbol{\sigma}}$ 是我们在 Weyl 场中已经非常熟悉的 $2 \times 2$ 矩阵。

对于纯空间旋转（$\boldsymbol{\phi} = 0$，旋转角 $\boldsymbol{\theta}$），两个子表示退化为相同的 $SU(2)$ 旋转矩阵 $e^{i\frac{1}{2}\boldsymbol{\theta}\cdot\boldsymbol{\sigma}}$，因此旋转矩阵不是块对角而是简单地在对角线上放置两个相同的矩阵：
\begin{equation}
    \mathscr{D}^{(1/2)}[R] = \begin{pmatrix} e^{i\frac{1}{2}\boldsymbol{\theta}\cdot\boldsymbol{\sigma}} & 0 \\ 0 & e^{i\frac{1}{2}\boldsymbol{\theta}\cdot\boldsymbol{\sigma}} \end{pmatrix}
\end{equation}
但对于Boost（$\boldsymbol{\theta} = 0$，快度 $\boldsymbol{\phi}$），两个子表示的矩阵符号相反：
\begin{equation}
    \mathscr{D}^{(1/2)}[L(p)] = \begin{pmatrix} e^{+\frac{1}{2}\boldsymbol{\phi}\cdot\boldsymbol{\sigma}} & 0 \\ 0 & e^{-\frac{1}{2}\boldsymbol{\phi}\cdot\boldsymbol{\sigma}} \end{pmatrix}
\end{equation}
正是 Boost 部分的这种"反对称性"，使得直和表示在宇称变换下能够保持封闭——$\gamma^0$ 矩阵交换上下两块，恰好补偿了 $\boldsymbol{\phi} \to -\boldsymbol{\phi}$ 带来的符号变化。

\paragraph{质量耦合与狄拉克方程}
对于无质量的粒子，左手和右手分量可以独立存在，分别满足 Weyl 方程。在动量空间中：
\begin{align}
    p_\mu \bar{\sigma}^\mu \psi_R &= (p^0 + \boldsymbol{p}\cdot\boldsymbol{\sigma}) \psi_R = 0 \\
    p_\mu \sigma^\mu \psi_L &= (p^0 - \boldsymbol{p}\cdot\boldsymbol{\sigma}) \psi_L = 0
\end{align}
其中 $\sigma^\mu = (I, \boldsymbol{\sigma})$，$\bar{\sigma}^\mu = (I, -\boldsymbol{\sigma})$。
当粒子具有静止质量 $m \neq 0$ 时，在静止参考系中（$\boldsymbol{p}=0, p^0=m$），洛伦兹群退化为空间旋转群 $SU(2)$。此时，左手 Weyl 旋量与右手 Weyl 旋量在纯旋转下变换完全相同，它们无法通过旋转来区分。因此，一个有质量的粒子在静止时不具有确定的手征性——左手态和右手态必须不可避免地通过质量 $m$ 发生耦合。

满足洛伦兹协变性且量纲正确的唯一线性交叉耦合方式为：
\begin{align}
    p_\mu \bar{\sigma}^\mu \psi_R &= m \psi_L \label{eq:weyl_coupled1} \\
    p_\mu \sigma^\mu \psi_L &= m \psi_R \label{eq:weyl_coupled2}
\end{align}
物理图像极为深刻： 这两个方程告诉我们，一个有质量的费米子在真空中传播时，实际上是在左手态和右手态之间不断地来回振荡。质量 $m$ 在这里充当的是"翻转手征性的耦合强度"。当 $m = 0$ 时，两个方程退耦，手征性成为守恒量，回到独立的无质量 Weyl 方程。
为了将上述两个 $2 \times 2$ 的耦合方程写成一个统一的 $4 \times 4$ 矩阵方程，我们引入分块矩阵结构：
\begin{equation}
    \begin{pmatrix} 0 & p_\mu \bar{\sigma}^\mu \\ p_\mu \sigma^\mu & 0 \end{pmatrix} \begin{pmatrix} \psi_R \\ \psi_L \end{pmatrix} = m \begin{pmatrix} \psi_R \\ \psi_L \end{pmatrix}
\end{equation}
左侧的 $4 \times 4$ 矩阵恰好可以分解为动量 $p_\mu$ 乘以一组 $4 \times 4$ 矩阵。我们定义这组矩阵为 Dirac $\gamma$ 矩阵。
\topictitle{Clifford 代数}
在手征表示下，为了让上述分块矩阵精确地写成 $\gamma^\mu p_\mu$ 的形式，我们定义：
\begin{equation}
    \gamma^0 = \begin{pmatrix} 0 & I \\ I & 0 \end{pmatrix}, \quad \gamma^i = \begin{pmatrix} 0 & -\sigma^i \\ \sigma^i & 0 \end{pmatrix}
\end{equation}
容易验证：
\begin{equation}
    \gamma^\mu p_\mu = \gamma^0 p^0 + \gamma^i p_i = \begin{pmatrix} 0 & p^0 I - \boldsymbol{p}\cdot\boldsymbol{\sigma} \\ p^0 I + \boldsymbol{p}\cdot\boldsymbol{\sigma} & 0 \end{pmatrix} = \begin{pmatrix} 0 & p_\mu \bar{\sigma}^\mu \\ p_\mu \sigma^\mu & 0 \end{pmatrix}
\end{equation}
（其中 $p_i = -p^i$，因此 $\gamma^i p_i = -\gamma^i p^i$。）
这恰好重现了我们的耦合方程！于是，统一的波动方程可以写为极为紧凑的狄拉克方程：
\begin{equation}
    (\gamma^\mu p_\mu - m)\psi = 0 \quad \xrightarrow{\text{坐标空间}} \quad (i\gamma^\mu \partial_\mu - m)\psi(x) = 0
\end{equation}
现在我们来推导 $\gamma$ 矩阵满足的代数关系。通过直接的 $2 \times 2$ 分块矩阵乘法：
\vspace{2ex}\par 
\textbf{计算 $\gamma^0 \gamma^0$：}
\begin{equation}
    \gamma^0 \gamma^0 = \begin{pmatrix} 0 & I \\ I & 0 \end{pmatrix} \begin{pmatrix} 0 & I \\ I & 0 \end{pmatrix} = \begin{pmatrix} I & 0 \\ 0 & I \end{pmatrix} = I_{4\times 4}
\end{equation}

\textbf{计算 $\gamma^i \gamma^j + \gamma^j \gamma^i$：}
\begin{equation}
    \gamma^i \gamma^j = \begin{pmatrix} 0 & -\sigma^i \\ \sigma^i & 0 \end{pmatrix} \begin{pmatrix} 0 & -\sigma^j \\ \sigma^j & 0 \end{pmatrix} = \begin{pmatrix} -\sigma^i\sigma^j & 0 \\ 0 & -\sigma^i\sigma^j \end{pmatrix}
\end{equation}
因此：
\begin{equation}
    \gamma^i \gamma^j + \gamma^j \gamma^i = \begin{pmatrix} -\{\sigma^i, \sigma^j\} & 0 \\ 0 & -\{\sigma^i, \sigma^j\} \end{pmatrix}
\end{equation}
由于泡利矩阵满足 $\{\sigma^i, \sigma^j\} = 2\delta^{ij} I$，我们得到 $\gamma^i \gamma^j + \gamma^j \gamma^i = -2\delta^{ij} I_{4\times 4}$。

\textbf{计算 $\gamma^0 \gamma^i + \gamma^i \gamma^0$：}
\begin{equation}
    \gamma^0 \gamma^i = \begin{pmatrix} 0 & I \\ I & 0 \end{pmatrix} \begin{pmatrix} 0 & -\sigma^i \\ \sigma^i & 0 \end{pmatrix} = \begin{pmatrix} \sigma^i & 0 \\ 0 & -\sigma^i \end{pmatrix}
\end{equation}
\begin{equation}
    \gamma^i \gamma^0 = \begin{pmatrix} 0 & -\sigma^i \\ \sigma^i & 0 \end{pmatrix} \begin{pmatrix} 0 & I \\ I & 0 \end{pmatrix} = \begin{pmatrix} -\sigma^i & 0 \\ 0 & \sigma^i \end{pmatrix}
\end{equation}
因此 $\gamma^0 \gamma^i + \gamma^i \gamma^0 = 0$。
\vspace{2ex}\par 
综合以上三种情形，我们得到了极其重要的 Clifford 代数：
\begin{equation}
    \boxed{\{\gamma^\mu, \gamma^\nu\} = \gamma^\mu \gamma^\nu + \gamma^\nu \gamma^\mu = 2\eta^{\mu\nu} I_{4\times 4}} \label{eq:clifford}
\end{equation}
这个代数关系保证了狄拉克方程的解自动满足 Klein-Gordon 方程（即正确的能量-动量关系）：
\begin{equation}
    (\gamma^\mu p_\mu + m)(\gamma^\nu p_\nu - m)\psi = (\gamma^\mu \gamma^\nu p_\mu p_\nu - m^2)\psi = \left(\frac{1}{2}\{\gamma^\mu,\gamma^\nu\}p_\mu p_\nu - m^2\right)\psi = (p^2 - m^2)\psi = 0
\end{equation}
为了方便地提取四分量旋量中的左手和右手分量，我们定义第五个矩阵 $\gamma^5$：
\begin{equation}
    \gamma^5 \equiv i\gamma^0\gamma^1\gamma^2\gamma^3 = \begin{pmatrix} -I & 0 \\ 0 & I \end{pmatrix}
\end{equation}
它满足如下代数关系：
\begin{equation}
    (\gamma^5)^2 = I, \quad \{\gamma^5, \gamma^\mu\} = 0 \quad (\text{即} \ \gamma^5 \gamma^\mu = -\gamma^\mu \gamma^5)
\end{equation}
作用在四分量旋量上时：
\begin{equation}
    \gamma^5 \psi = \begin{pmatrix} -I & 0 \\ 0 & I \end{pmatrix} \begin{pmatrix} \psi_R \\ \psi_L \end{pmatrix} = \begin{pmatrix} -\psi_R \\ \psi_L \end{pmatrix}
\end{equation}
利用 $\gamma^5$ 的本征值 $\pm 1$，可以定义精准的手征投影算符：
\begin{equation}
    P_L = \frac{1 + \gamma^5}{2} = \begin{pmatrix} 0 & 0 \\ 0 & I \end{pmatrix}, \quad P_R = \frac{1 - \gamma^5}{2} = \begin{pmatrix} I & 0 \\ 0 & 0 \end{pmatrix}
\end{equation}
满足 $P_L + P_R = I$，$P_L P_R = 0$，$P_L^2 = P_L$。
\vspace{2ex}\par 
在洛伦兹变换 $x \to \Lambda x$ 下，四分量旋量的变换为 $\psi'(x') = S(\Lambda) \psi(x)$。
对于无穷小变换 $\Lambda^\mu_{\,\,\nu} \approx \delta^\mu_{\,\,\nu} + \omega^\mu_{\,\,\nu}$，旋量空间的变换矩阵可以由 Clifford 代数直接构造：
\begin{equation}
    S(1 + \omega) = I - \frac{i}{4} \omega_{\mu\nu} \Sigma^{\mu\nu}, \quad \Sigma^{\mu\nu} = \frac{i}{2}[\gamma^\mu, \gamma^\nu]
\end{equation}

\textbf{致命问题：洛伦兹群是非紧致群。}
对于空间旋转（$\omega_{ij}$），$\Sigma^{ij}$ 是厄米矩阵，因此旋转部分的 $S$ 是幺正的（$S^\dagger = S^{-1}$）。
但对于 Boost（$\omega_{0i}$），$\Sigma^{0i}$ 是反厄米矩阵，Boost 部分的 $S$ 不是幺正的（$S^\dagger \neq S^{-1}$）。

这直接导致了如下灾难：非相对论量子力学中熟悉的概率密度 $\psi^\dagger \psi$ 在 Boost 下不再是不变量！
\begin{equation}
    \psi'^\dagger \psi' = \psi^\dagger S^\dagger S \psi \neq \psi^\dagger \psi \quad (\text{因为} \ S^\dagger S \neq I)
\end{equation}

为了构造真正的洛伦兹标量，我们需要一个矩阵来"修正" $S^\dagger$ 的非幺正性。利用 $\gamma$ 矩阵的显式形式，可以验证：
\begin{equation}
    (\gamma^0)^\dagger = \gamma^0, \quad (\gamma^i)^\dagger = -\gamma^i \quad \Longrightarrow \quad \gamma^0 (\gamma^\mu)^\dagger \gamma^0 = \gamma^\mu
\end{equation}
这导致了关键恒等式 $\gamma^0 S^\dagger \gamma^0 = S^{-1}$，即 $S^\dagger = \gamma^0 S^{-1} \gamma^0$。

因此，我们定义狄拉克共轭 (Dirac Adjoint)：
\begin{equation}
    \bar{\psi} \equiv \psi^\dagger \gamma^0
\end{equation}
它在洛伦兹变换下的行为极其优美：
\begin{equation}
    \bar{\psi}' = \psi'^\dagger \gamma^0 = \psi^\dagger S^\dagger \gamma^0 = \psi^\dagger (\gamma^0 S^{-1} \gamma^0) \gamma^0 = \psi^\dagger \gamma^0 S^{-1} = \bar{\psi} S^{-1}
\end{equation}
从而保证了 $\bar{\psi}\psi$ 是真正的洛伦兹标量：
\begin{equation}
    \bar{\psi}'\psi' = \bar{\psi} S^{-1} S \psi = \bar{\psi}\psi
\end{equation}
这也是为什么拉格朗日量中的质量项必须写成 $m\bar{\psi}\psi$ 的根本原因。
\vspace{2ex}\par 
任意 $4 \times 4$ 复矩阵有 16 个独立分量，恰好可以用 16 个线性独立的 $\gamma$ 矩阵组合作为基底展开。将这些基底夹在 $\bar{\psi}$ 和 $\psi$ 之间，可以构造出 5 种具有确定洛伦兹张量变换性质的双线性形式：

\begin{center}
\begin{tabular}{|c|c|c|c|}
\hline
\textbf{类型} & \textbf{基底矩阵} & \textbf{自由度数} & \textbf{变换性质} \\
\hline
标量 (S) & $I$ & 1 & $\bar{\psi}\psi \to \bar{\psi}\psi$ \\
\hline
赝标量 (P) & $i\gamma^5$ & 1 & $\bar{\psi}i\gamma^5\psi \xrightarrow{\mathcal{P}} -\bar{\psi}i\gamma^5\psi$ \\
\hline
矢量 (V) & $\gamma^\mu$ & 4 & $\bar{\psi}\gamma^\mu\psi \to \Lambda^\mu_{\,\,\nu}\bar{\psi}\gamma^\nu\psi$ \\
\hline
轴矢量 (A) & $\gamma^\mu\gamma^5$ & 4 & $\bar{\psi}\gamma^\mu\gamma^5\psi \xrightarrow{\mathcal{P}} -\Lambda^\mu_{\,\,\nu}\bar{\psi}\gamma^\nu\gamma^5\psi$ \\
\hline
张量 (T) & $\Sigma^{\mu\nu} = \frac{i}{2}[\gamma^\mu,\gamma^\nu]$ & 6 & $\bar{\psi}\Sigma^{\mu\nu}\psi \to \Lambda^\mu_{\,\,\rho}\Lambda^\nu_{\,\,\sigma}\bar{\psi}\Sigma^{\rho\sigma}\psi$ \\
\hline
\end{tabular}
\end{center}
总计 $1 + 1 + 4 + 4 + 6 = 16$ 个，完美匹配 $4 \times 4$ 矩阵空间的维度。
\topictitle{分立对称性}
在当前的经典容器层面，分立对称性表现为旋量空间中的特殊矩阵：

\textbf{1. 空间反演 }

空间反演使坐标 $\boldsymbol{x} \to -\boldsymbol{x}$，在狄拉克方程中翻转空间微商 $\partial_i \to -\partial_i$。我们需要一个矩阵能够同时实现"交换左右手分量"和"翻转空间 $\gamma^i$ 的符号"。根据 Clifford 代数 \eqref{eq:clifford}：
\begin{equation}
    \gamma^0 \gamma^i (\gamma^0)^{-1} = \gamma^0 \gamma^i \gamma^0 = -\gamma^i, \quad \gamma^0 \gamma^0 (\gamma^0)^{-1} = \gamma^0
\end{equation}
因此宇称矩阵正是 $\gamma^0$ 本身：
\begin{equation}
    \psi_P(t, \boldsymbol{x}) = \eta_P \gamma^0 \psi(t, -\boldsymbol{x})
\end{equation}

\textbf{2. 电荷共轭 }

电荷共轭将粒子解映射为反粒子解。对狄拉克方程 $(i\gamma^\mu\partial_\mu - m)\psi = 0$ 取复共轭后变为 $(-i\gamma^{\mu *}\partial_\mu - m)\psi^* = 0$。我们需要一个矩阵 $C$ 将 $-\gamma^{\mu *}$ 变回 $+\gamma^\mu$，即 $C^{-1}\gamma^\mu C = -(\gamma^\mu)^T$。在手征表示下：
\begin{equation}
    C = i\gamma^2\gamma^0, \quad \psi_C(t, \boldsymbol{x}) = \eta_C C \bar{\psi}^T(t, \boldsymbol{x})
\end{equation}

\textbf{3. 时间反演}

时间反演是反线性算符，包含复共轭和时间坐标翻转 $t \to -t$：
\begin{equation}
    \psi_T(t, \boldsymbol{x}) = \eta_T B \psi^*(-t, \boldsymbol{x}), \quad \text{其中} \ B^{-1}\gamma^\mu B = (\gamma^\mu)^*
\end{equation}
\subsection{因果狄拉克场 }
一般形式的场是：
\begin{equation}
    \Psi_{ab}(x) =(\frac{1}{2\pi})^\frac{3}{2}\sum_{\sigma}\int d^3p \, \left\{(-)^{2B} u_{ab}(\boldsymbol{q},\sigma)b^\dagger(\boldsymbol{q},\sigma) e^{-ipx} + v_{ab}(\boldsymbol{q},\sigma) a(\boldsymbol{q},\sigma)e^{ipx} \right\}
\end{equation}
上述我们知道静止系系数 $u(0)$ 和 $v(0)$ 是正比于Clebsch-Gordan  系数 
\begin{equation}
    u_{ab}(0, \sigma) = \frac{1}{\sqrt{2m}} \langle A, a; B, b | j, \sigma \rangle
    \label{eq:u0_general_1/2}
\end{equation}
通过 Boost 获得有限动量系数 $u(\boldsymbol{q})$ 和 $v(\boldsymbol{q})$的一般公式是：
\begin{equation}
    u_{ab}(\boldsymbol{q}, \sigma) = \frac{1}{\sqrt{2q^0}} \sum_{\bar{a},\bar{b}} \left[e^{-\boldsymbol{\phi}\cdot\mathbf{J}^{(A)}}\right]_{a\bar{a}} \left[e^{+\boldsymbol{\phi}\cdot\mathbf{J}^{(B)}}\right]_{b\bar{b}} \langle A,\bar{a}; B,\bar{b} | j, \sigma \rangle
    \label{eq:u_finite_general_1/2}
\end{equation}
其中$\boldsymbol{\phi} = \hat{\boldsymbol{q}} \text{arctanh}(|\boldsymbol{q}|/q^0)$
\vspace{0.5cm}
\begin{center}
    \large\text{$(A, B) = (1/2, 0)\oplus(1/2, 0)$}
\end{center}
\vspace{0.5cm}
对于自旋 $j=1/2$ 的有质量粒子，在静止参考系（$\boldsymbol{p}=\mathbf{0}$，$p^0=m$）中，狄拉克方程 $(\gamma^\mu p_\mu - m)u=0$ 退化为：
\begin{equation}
    (m\gamma^0 - m)u(0,\sigma) = 0 \quad\Longrightarrow\quad \gamma^0 u(0,\sigma) = u(0,\sigma)
\end{equation}
这是 $\gamma^0$ 矩阵的本征值为 $+1$ 的本征方程。在手征表示下 $\gamma^0 = \begin{pmatrix} 0&I\\I&0 \end{pmatrix}$，写出分量形式：
\begin{equation}
    \begin{pmatrix} 0&I\\I&0 \end{pmatrix}\begin{pmatrix} u_R \\ u_L \end{pmatrix} = \begin{pmatrix} u_R \\ u_L \end{pmatrix} \quad\Longrightarrow\quad u_L = u_R
\end{equation}
因此本征值 $+1$ 的本征空间是二维的（$u_R = u_L$ 可取两个线性独立基矢），对应两个自旋投影 $\sigma = \pm 1/2$。

选择洛伦兹标量归一化 $\bar{u}u = u^\dagger\gamma^0 u = 2m$，逐一写出具体的四分量列向量：

\textbf{$\sigma = +1/2$：}
\begin{equation}
    u(0, +\tfrac{1}{2}) = \sqrt{m}\begin{pmatrix} 1\\0\\1\\0 \end{pmatrix} = \sqrt{m}\begin{pmatrix} \xi_{+1/2} \\ \xi_{+1/2} \end{pmatrix}
\end{equation}

\textbf{$\sigma = -1/2$：}
\begin{equation}
    u(0, -\tfrac{1}{2}) = \sqrt{m}\begin{pmatrix} 0\\1\\0\\1 \end{pmatrix} = \sqrt{m}\begin{pmatrix} \xi_{-1/2} \\ \xi_{-1/2} \end{pmatrix}
\end{equation}
其中 $\xi_{+1/2} = \binom{1}{0}$，$\xi_{-1/2} = \binom{0}{1}$。简记为 $u(0,\sigma) = \sqrt{m}\binom{\xi_\sigma}{\xi_\sigma}$。

\textbf{$\sigma = +1/2$：} 
\begin{equation}
    (-)^{1}(-)^{1+1/2}u(0, +\tfrac{1}{2}) = \sqrt{m}\begin{pmatrix} 0\\1\\0\\-1 \end{pmatrix} = \sqrt{m}\begin{pmatrix} \chi_{+1/2} \\ -\chi_{+1/2} \end{pmatrix}
\end{equation}

\textbf{$\sigma = -1/2$：} 
\begin{equation}
        (-)^{1}(-)^{1+1/2}u(0, -\tfrac{1}{2})= \sqrt{m}\begin{pmatrix} -1\\0\\1\\0 \end{pmatrix} = \sqrt{m}\begin{pmatrix} \chi_{-1/2} \\ -\chi_{-1/2} \end{pmatrix}
\end{equation}
简记为 $v(0,\sigma) = \sqrt{m}\binom{\chi_\sigma}{-\chi_\sigma}$。
\vspace{2ex}\par
在直和表示 $(1/2,0)\oplus(0,1/2)$ 中，Boost 矩阵是对角分块的（上块按右手加速，下块按左手加速）：
\begin{equation}
    S(L(p)) = \begin{pmatrix} e^{+\frac{1}{2}\boldsymbol{\phi}\cdot\boldsymbol{\sigma}} & 0 \\ 0 & e^{-\frac{1}{2}\boldsymbol{\phi}\cdot\boldsymbol{\sigma}} \end{pmatrix}
\end{equation}
利用泡利矩阵的性质 $(\hat{\boldsymbol{p}}\cdot\boldsymbol{\sigma})^2 = I$ 以及双曲半角公式（$\cosh\phi = p^0/m$，$\sinh\phi = |\boldsymbol{p}|/m$）：
\begin{align}
    e^{+\frac{1}{2}\boldsymbol{\phi}\cdot\boldsymbol{\sigma}} &= \cosh\frac{\phi}{2}\cdot I + \sinh\frac{\phi}{2}\cdot\hat{\boldsymbol{p}}\cdot\boldsymbol{\sigma} = \frac{(p^0+m)I + \boldsymbol{p}\cdot\boldsymbol{\sigma}}{\sqrt{2m(p^0+m)}} \\
    e^{-\frac{1}{2}\boldsymbol{\phi}\cdot\boldsymbol{\sigma}} &= \cosh\frac{\phi}{2}\cdot I - \sinh\frac{\phi}{2}\cdot\hat{\boldsymbol{p}}\cdot\boldsymbol{\sigma} = \frac{(p^0+m)I - \boldsymbol{p}\cdot\boldsymbol{\sigma}}{\sqrt{2m(p^0+m)}}
\end{align}
将 Boost 矩阵作用于静止系旋量 $u(0,\sigma) = \sqrt{m}\binom{\xi_\sigma}{\xi_\sigma}$：
\begin{equation}
    u(\boldsymbol{p},\sigma) = S(L(p))\, u(0,\sigma) = \sqrt{m}\begin{pmatrix} e^{+\frac{1}{2}\boldsymbol{\phi}\cdot\boldsymbol{\sigma}}\xi_\sigma \\ e^{-\frac{1}{2}\boldsymbol{\phi}\cdot\boldsymbol{\sigma}}\xi_\sigma \end{pmatrix}
\end{equation}
四维记号下： $\sqrt{p\cdot\bar{\sigma}} \equiv e^{+\frac{1}{2}\boldsymbol{\phi}\cdot\boldsymbol{\sigma}}\sqrt{m}$ 和 $\sqrt{p\cdot\sigma} \equiv e^{-\frac{1}{2}\boldsymbol{\phi}\cdot\boldsymbol{\sigma}}\sqrt{m}$（其中 $p\cdot\sigma = p^0 I - \boldsymbol{p}\cdot\boldsymbol{\sigma}$，$p\cdot\bar{\sigma} = p^0 I + \boldsymbol{p}\cdot\boldsymbol{\sigma}$），则：
\begin{equation}
    \boxed{u(\boldsymbol{p},\sigma) = \begin{pmatrix} \sqrt{p\cdot\bar{\sigma}}\,\xi_\sigma \\ \sqrt{p\cdot\sigma}\,\xi_\sigma \end{pmatrix}}
\end{equation}
将 Boost 矩阵作用于 $v(0,\sigma) = \sqrt{m}\binom{\chi_\sigma}{-\chi_\sigma}$（注意下半部分的负号在 Boost 后被保留）：
\begin{equation}
    \boxed{v(\boldsymbol{p},\sigma) = \begin{pmatrix} \sqrt{p\cdot\bar{\sigma}}\,\chi_\sigma \\ -\sqrt{p\cdot\sigma}\,\chi_\sigma \end{pmatrix}}
\end{equation}
这些波函数在任意惯性系下自动满足 $(\slashed{p}-m)u=0$ 和 $(\slashed{p}+m)v=0$。
\vspace{2ex}\par
为了得到洛伦兹协变的结果，我们计算 $u\bar{u}$（其中 $\bar{u}=u^\dagger\gamma^0$）而非 $uu^\dagger$。

对 $\sigma = \pm 1/2$ 求和，利用 $\sum_\sigma\xi_\sigma\xi_\sigma^\dagger = I$，逐块计算 $4\times 4$ 矩阵的四个 $2\times 2$ 子块：
\begin{itemize}
    \item 左上块：$\sqrt{p\cdot\bar{\sigma}}\cdot I\cdot\sqrt{p\cdot\sigma} = \sqrt{(p\cdot\bar{\sigma})(p\cdot\sigma)}$。利用 $(p\cdot\bar{\sigma})(p\cdot\sigma) = (p^0)^2 - |\boldsymbol{p}|^2 = m^2$，得 $= mI$。
    \item 右上块：$\sqrt{p\cdot\bar{\sigma}}\cdot I\cdot\sqrt{p\cdot\bar{\sigma}} = p\cdot\bar{\sigma} = p^0 I + \boldsymbol{p}\cdot\boldsymbol{\sigma}$。
    \item 左下块：$\sqrt{p\cdot\sigma}\cdot I\cdot\sqrt{p\cdot\sigma} = p\cdot\sigma = p^0 I - \boldsymbol{p}\cdot\boldsymbol{\sigma}$。
    \item 右下块：$\sqrt{p\cdot\sigma}\cdot I\cdot\sqrt{p\cdot\bar{\sigma}} = mI$。
\end{itemize}
因此：
\begin{equation}
    \sum_\sigma u(\boldsymbol{p},\sigma)\bar{u}(\boldsymbol{p},\sigma) = \begin{pmatrix} mI & p\cdot\bar{\sigma} \\ p\cdot\sigma & mI \end{pmatrix} = \gamma^\mu p_\mu + m \equiv \slashed{p}+m
    \label{eq:pol_u}
\end{equation}
对于反粒子旋量 $v$，由于下半部分负号导致对角块反号：
\begin{equation}
   \sum_\sigma v(\boldsymbol{p},\sigma)\bar{v}(\boldsymbol{p},\sigma) = \slashed{p}-m
    \label{eq:pol_v}
\end{equation}



\subsection{因果矢量场 }
矢量场对应于洛伦兹群的 $(1/2,1/2)$ 表示，描述自旋 $j=1$ 的有质量粒子。本节我们将详细展示其静止系系数如何从角动量耦合中被唯一确定，以及如何通过 Boost 获得有限动量的系数。
一般形式的场是：
\begin{equation}
    \Psi_{ab}(x) =(\frac{1}{2\pi})^\frac{3}{2}\sum_{\sigma}\int d^3p \, \left\{ u_{ab}(\boldsymbol{q},\sigma)b^\dagger(\boldsymbol{q},\sigma) e^{-ipx} + (-)^{2B}(-)^{j+\sigma}u_{ab}(\boldsymbol{q},\sigma) a(\boldsymbol{q},\sigma)e^{ipx} \right\}
\end{equation}
上述我们知道静止系系数 $u(0)$ 和 $v(0)$ 是正比于Clebsch-Gordan  系数 
\begin{equation}
    u_{ab}(0, \sigma) = \frac{1}{\sqrt{2m}} \langle A, a; B, b | j, \sigma \rangle
    \label{eq:u0_general_1/2}
\end{equation}
通过 Boost 获得有限动量系数 $u(\boldsymbol{q})$ 和 $v(\boldsymbol{q})$的一般公式是：
\begin{equation}
    u_{ab}(\boldsymbol{q}, \sigma) = \frac{1}{\sqrt{2q^0}} \sum_{\bar{a},\bar{b}} \left[e^{-\boldsymbol{\phi}\cdot\mathbf{J}^{(A)}}\right]_{a\bar{a}} \left[e^{+\boldsymbol{\phi}\cdot\mathbf{J}^{(B)}}\right]_{b\bar{b}} \langle A,\bar{a}; B,\bar{b} | j, \sigma \rangle
    \label{eq:u_finite_general_1/2}
\end{equation}
其中$\boldsymbol{\phi} = \hat{\boldsymbol{q}} \text{arctanh}(|\boldsymbol{q}|/q^0)$
\vspace{0.5cm}
\begin{center}
    \large\text{$(A, B) = (1/2, 1/2)=(1, 0)\oplus(0,0)$}
\end{center}
\vspace{0.5cm}
\paragraph{角动量加法}
对于 $(A,B)=(1/2,1/2)$ 表示，场的指标为 $a=\pm 1/2$（对应 $A=1/2$）和 $b=\pm 1/2$（对应 $B=1/2$）。表示空间的维度为 $(2A+1)(2B+1)=2\times 2=4$。

在静止参考系（$\mathbf{p}=\mathbf{0}$）中，洛伦兹群退化为旋转群 $SU(2)$。总角动量 $\mathbf{J}=\mathbf{J}^{(A)}+\mathbf{J}^{(B)}$ 是两个自旋 $1/2$ 的耦合。根据角动量加法定理：
\begin{equation}
    \frac{1}{2}\otimes\frac{1}{2} = 1\oplus 0
\end{equation}
即 $4$ 维空间分解为一个 $3$ 维的 $j=1$（三重态）和一个 $1$ 维的 $j=0$（单态）不可约表示。
\vspace{2ex}\par
我们逐一写出全部的 C-G 系数 $\langle 1/2,a;\,1/2,b\,|\,j,\sigma\rangle$。
\textbf{$j=1$ 三重态（3 个状态）：}
\begin{align}
    |1,+1\rangle &= |a{=}+\tfrac{1}{2},\, b{=}+\tfrac{1}{2}\rangle \\
    |1,\;0\;\rangle &= \frac{1}{\sqrt{2}}\Big(|a{=}+\tfrac{1}{2},\, b{=}-\tfrac{1}{2}\rangle + |a{=}-\tfrac{1}{2},\, b{=}+\tfrac{1}{2}\rangle\Big) \\
    |1,-1\rangle &= |a{=}-\tfrac{1}{2},\, b{=}-\tfrac{1}{2}\rangle
\end{align}
对应的非零 C-G 系数为：
\begin{equation}
    \begin{aligned}
        \langle +\tfrac{1}{2},+\tfrac{1}{2}\,|\,1,+1\rangle &= 1 \\
        \langle +\tfrac{1}{2},-\tfrac{1}{2}\,|\,1,\;0\;\rangle &= \tfrac{1}{\sqrt{2}},\quad \langle -\tfrac{1}{2},+\tfrac{1}{2}\,|\,1,\;0\;\rangle = \tfrac{1}{\sqrt{2}} \\
        \langle -\tfrac{1}{2},-\tfrac{1}{2}\,|\,1,-1\rangle &= 1
    \end{aligned}
\end{equation}

\textbf{$j=0$ 单态（1 个状态）：}
\begin{equation}
    |0,0\rangle = \frac{1}{\sqrt{2}}\Big(|a{=}+\tfrac{1}{2},\, b{=}-\tfrac{1}{2}\rangle - |a{=}-\tfrac{1}{2},\, b{=}+\tfrac{1}{2}\rangle\Big)
\end{equation}
对应非零 C-G 系数为：
\begin{equation}
    \langle +\tfrac{1}{2},-\tfrac{1}{2}\,|\,0,0\rangle = +\tfrac{1}{\sqrt{2}},\quad \langle -\tfrac{1}{2},+\tfrac{1}{2}\,|\,0,0\rangle = -\tfrac{1}{\sqrt{2}}
\end{equation}
\begin{itemize}
    \item $j=0$ 部分：时间分量(标量):$j=0$ 态在旋转下不变（标量）。在四维矢量 $A^\mu$ 中，哪个分量在空间旋转下不变？答案是时间分量 $A^0$。

    利用泡利矩阵 $\sigma^\mu=(I,\boldsymbol{\sigma})$ 将四维矢量指标 $\mu$ 与旋量指标 $(a,b)$ 联系起来，$(1/2,1/2)$ 表示空间中的 $2\times 2$ 矩阵可以展开为：
    \begin{equation}
    \phi_{ab} = (\sigma^\mu)_{ab}\, A_\mu = A_0\, I_{ab} + A_i\, (\sigma^i)_{ab}
    \end{equation}
    时间分量 $A_0$ 对应 $\sigma^0=I$（$2\times 2$ 单位阵）。而单位阵 $I_{ab}\propto\delta_{ab}$ 恰好对应 $j=0$ 单态的结构——$j=0$ 的 C-G 系数给出的是反对称组合 $\propto|{+}{-}\rangle-|{-}{+}\rangle$，在缩并为 $2\times 2$ 矩阵时正是正比于 $\epsilon_{ab}$（Levi-Civita 符号），通过指标升降后等价于 $\delta_{ab}$。所以$j=0$ 部分精确对应于四维矢量的时间分量 $A^0$（旋转标量）。
    \item$j=1$ 部分：空间分量(矢量):$j=1$ 态在旋转下按自旋 $1$ 变换。在四维矢量中，空间分量 $(A^1,A^2,A^3)$ 在旋转下构成三维矢量。

    空间分量 $A_i$ 对应泡利矩阵 $\sigma^i$。而 $\sigma^1,\sigma^2,\sigma^3$ 构成的无迹 $2\times 2$ 矩阵空间，恰好对应 $j=1$ 三重态——$j=1$ 的 C-G 系数（对称组合 $|{+}{+}\rangle$，$|{+}{-}\rangle+|{-}{+}\rangle$，$|{-}{-}\rangle$）在缩并为矩阵时给出的正是泡利矩阵的结构。$j=1$ 部分精确对应于四维矢量的空间分量 $(A^1,A^2,A^3)$（旋转矢量）。
\end{itemize}


$j=1$ 态在旋转下按自旋 $1$ 变换。在四维矢量中，空间分量 $(A^1,A^2,A^3)$ 在旋转下构成三维矢量。

空间分量 $A_i$ 对应泡利矩阵 $\sigma^i$。而 $\sigma^1,\sigma^2,\sigma^3$ 构成的无迹 $2\times 2$ 矩阵空间，恰好对应 $j=1$ 三重态——$j=1$ 的 C-G 系数（对称组合 $|{+}{+}\rangle$，$|{+}{-}\rangle+|{-}{+}\rangle$，$|{-}{-}\rangle$）在缩并为矩阵时给出的正是泡利矩阵的结构。

\textbf{结论：$j=1$ 部分精确对应于四维矢量的空间分量 $(A^1,A^2,A^3)$（旋转矢量）。}

\paragraph{$j=0$部分与横向条件}

我们要描述的是自旋 $j=1$ 的粒子（而非自旋 $0$）。因此在构造静止系系数时，只能使用 $j=1$ 的 C-G 系数：
\begin{equation}
    u_{ab}(0,\sigma) = \frac{1}{\sqrt{2m}}\langle\tfrac{1}{2},a;\,\tfrac{1}{2},b\,|\,1,\sigma\rangle
\end{equation}
而\textbf{绝不使用} $j=0$ 的 C-G 系数。在四维矢量语言中，这精确等价于要求静止系的极化矢量没有时间分量：
\begin{equation}
    \varepsilon^0(0,\sigma)=0 \quad\Longrightarrow\quad p_\mu\varepsilon^\mu(0,\sigma)=m\cdot\varepsilon^0=0
\end{equation}
这就是横向条件 $p_\mu\varepsilon^\mu=0$ 的群论起源——它并非外加的人为约束，而是我们只取 $(1/2,1/2)$ 直积表示中 $j=1$ 不可约部分的自然后果。
\vspace{2ex}\par
排除了 $j=0$（时间分量）后，极化矢量完全由空间自旋 $1$ 的球面基矢给出：
\begin{align}
    u^\mu(0,+1) &= \varepsilon^\mu(0,+1) = \frac{1}{\sqrt{2}}(0,-1,-i,0) \\
    u^\mu(0,\;0\;) &= \varepsilon^\mu(0,\;0\;) = (0,0,0,1) \\
    u^\mu(0,-1) &= \varepsilon^\mu(0,-1) = \frac{1}{\sqrt{2}}(0,1,-i,0)
\end{align}
满足正交归一关系 $\varepsilon^*_\mu(0,\sigma)\varepsilon^\mu(0,\sigma')=-\delta_{\sigma\sigma'}$。反粒子系数根据一般规律：$v^\mu(0,\sigma)=(-1)^{1-\sigma}u^\mu(0,-\sigma)=\varepsilon^{\mu*}(0,\sigma)$。
\vspace{2ex}\par

由于 $(1/2,1/2)$ 表示等价于四维矢量表示，Boost 矩阵就是经典的 $4\times 4$ 洛伦兹变换矩阵 $\Lambda^\mu_{\;\nu}(L(p))$。一般公式为：
\begin{equation}
    u^\mu(\mathbf{p},\sigma) = \varepsilon^\mu(\mathbf{p},\sigma) = \Lambda^\mu_{\;\nu}(L(p))\,\varepsilon^\nu(0,\sigma)
\end{equation}

\vspace{2ex}\par 设粒子沿 $z$ 轴加速到动量 $|\mathbf{p}|$。洛伦兹 Boost 矩阵为：
\begin{equation}
    \Lambda^\mu_{\;\nu}(L(p)) = \begin{pmatrix} p^0/m & 0 & 0 & |\mathbf{p}|/m \\ 0&1&0&0 \\ 0&0&1&0 \\ |\mathbf{p}|/m & 0 & 0 & p^0/m \end{pmatrix}
\end{equation}
横向极化只有 $x,y$ 分量，垂直于 Boost 方向，Boost 矩阵对其完全不起作用：
\begin{equation}
    \varepsilon^\mu(\mathbf{p},\pm 1) = \Lambda^\mu_{\;\nu}\,\varepsilon^\nu(0,\pm 1) = \varepsilon^\mu(0,\pm 1) = \frac{1}{\sqrt{2}}(0,\mp 1,-i,0)
\end{equation}
纵向极化只有 $z$ 分量，平行于 Boost 方向，时间和空间分量发生混合：
\begin{equation}
    \varepsilon^\mu(\mathbf{p},0) = \begin{pmatrix} p^0/m&0&0&|\mathbf{p}|/m\\0&1&0&0\\0&0&1&0\\|\mathbf{p}|/m&0&0&p^0/m \end{pmatrix}\begin{pmatrix}0\\0\\0\\1\end{pmatrix} = \begin{pmatrix}|\mathbf{p}|/m\\0\\0\\p^0/m\end{pmatrix} = \left(\frac{|\mathbf{p}|}{m},\,\frac{p^0}{m}\hat{\mathbf{p}}\right)
\end{equation}
静止系中纯空间的纵向极化在 Boost 后获得了非零时间分量 $\varepsilon^0=|\mathbf{p}|/m$。但横向条件仍然成立：
\begin{equation}
    p_\mu\varepsilon^\mu(\mathbf{p},0) = p^0\cdot\frac{|\mathbf{p}|}{m} - |\mathbf{p}|\cdot\frac{p^0}{m} = 0
\end{equation}
\vspace{2ex}\par
对三种极化状态 $\sigma=\pm 1,0$ 求和，利用洛伦兹张量的协变性写为：
\begin{equation}
    \boxed{\sum_{\sigma=\pm 1,0}\varepsilon^\mu(\mathbf{p},\sigma)\,\varepsilon^{\nu*}(\mathbf{p},\sigma) = -\eta^{\mu\nu}+\frac{p^\mu p^\nu}{m^2}}
    \label{eq:vec_pol}
\end{equation}
物理含义：$-\eta^{\mu\nu}$ 对应全部 $4$ 个自由度的恒等投影；$+p^\mu p^\nu/m^2$ 精确扣除沿 $p^\mu$ 方向的非物理 $j=0$ 自由度（$4-1=3$ 个物理自由度）。



\paragraph{无质量极限与规范不变性}

当 $m\to 0$ 时，纵向极化 $\varepsilon^\mu(\mathbf{p},0)=(|\mathbf{p}|/m,\,p^0\hat{\mathbf{p}}/m)$ 发散，极化求和与传播子中的 $p^\mu p^\nu/m^2$ 项出现奇异性。

Wigner 小群分析表明，无质量矢量粒子只有 $2$ 个横向螺旋度（$\lambda=\pm 1$），纵向自由度不物理。发散项总正比于 $p^\mu$。若矢量场仅与守恒流 $J_\mu$ 耦合（$\partial^\mu J_\mu=0$，即 $p^\mu J_\mu=0$），则所有奇异项严格消失。

而要求流守恒等价于要求理论具有\textbf{局部规范不变性}：$A^\mu(x)\to A^\mu(x)+\partial^\mu\alpha(x)$。

\textbf{终极结论：}微观因果性与洛伦兹不变性共同证明，自然界要容纳无质量自旋 $1$ 粒子，就必须以规范对称性作为其存在的数学保障。

   